% ============================================================
% BAB 1 - PENDAHULUAN DAN KONTEKS DOKUMEN
% ============================================================

\subsection{Tujuan Dokumen}

Dokumen ini mendefinisikan kebutuhan perangkat lunak untuk PasarKita sebagai
komponen \textit{Demand Generator} B2C dalam ekosistem \textit{microservices}
UMKM RPL 2.

SRS ini berfungsi sebagai \textit{baseline} bersama antara pemilik produk,
pengembang, penguji, dan pengelola sistem. Dokumen tidak dimaksudkan sebagai
uraian konseptual, melainkan sebagai spesifikasi yang menerjemahkan perilaku
aplikasi ke dalam kebutuhan yang dapat dibangun, diuji, dan dipelihara. Setiap
kebutuhan dirumuskan agar memiliki ruang lingkup yang jelas, dasar implementasi
yang dapat ditelusuri, dan kriteria penerimaan yang dapat diverifikasi.

Rujukan format yang digunakan adalah praktik SRS yang umum dipakai di lingkungan
pengembangan perangkat lunak: dokumen dimulai dari tujuan dan \textit{scope},
dilanjutkan dengan \textit{product perspective}, \textit{user classes},
\textit{constraints}, asumsi, kebutuhan fungsional, kebutuhan non-fungsional,
\textit{external interfaces}, \textit{data requirements}, \textit{business rules},
\textit{acceptance criteria}, dan \textit{traceability matrix}. Dengan struktur
tersebut, dokumen dapat dipakai sebagai kontrak kerja teknis tanpa bergantung
pada penjelasan lisan.

% ----
\subsection{Ruang Lingkup Produk}

PasarKita adalah platform \textit{marketplace} digital B2C yang memfasilitasi
transaksi produk UMKM antara \textit{seller} dan \textit{buyer}. Sistem berjalan
sebagai aplikasi web \textit{fullstack} dengan \textit{frontend} Next.js 16.2
(di-\textit{deploy} ke \texttt{pasarkita.vercel.app}) dan \textit{backend}
Express.js berbasis \textit{serverless} (di-\textit{deploy} ke
\texttt{pasarkita-api.vercel.app}), menggunakan Supabase PostgreSQL sebagai
\textit{database}.

Dalam arsitektur ekosistem, PasarKita berperan sebagai \textbf{Demand
Generator}---inisiator transaksi di sisi permintaan. Semua transaksi keuangan
diproses oleh SmartBank (Kelompok 1), semua komunikasi antar-\textit{service}
melewati API Gateway (Kelompok 7), dan pengiriman dikoordinasikan oleh LogistiKita
(Kelompok 5). PasarKita tidak mengambil alih fungsi pembayaran, \textit{routing},
atau logistik dari \textit{service partner} tersebut.

Sistem PasarKita bertanggung jawab pada siklus transaksi \textit{marketplace}:
registrasi/autentikasi \textit{user}, pengelolaan katalog produk UMKM, keranjang
belanja, proses \textit{checkout} dengan kalkulasi \textit{fee}, koordinasi
pembayaran ke SmartBank, \textit{trigger} pengiriman ke LogistiKita, manajemen
status \textit{order}, dan \textit{dashboard analytics} untuk \textit{superadmin}.

% ----
\subsection{Out of Scope}

Dokumen ini tidak menspesifikasikan logika pembayaran internal SmartBank,
mekanisme \textit{routing} JWT API Gateway, algoritma penentuan tarif pengiriman
LogistiKita, ataupun manajemen akun dari \textit{service partner} lain.

\begin{longtable}{|L{3.5cm}|L{3cm}|L{8.5cm}|}
    \caption{Area di Luar dan di Dalam Lingkup SRS}
    \label{tab:scope} \\
    \hline
    \textbf{Area} & \textbf{Status Scope} & \textbf{Rasional} \\
    \hline
    \endfirsthead
    \hline
    \textbf{Area} & \textbf{Status Scope} & \textbf{Rasional} \\
    \hline
    \endhead
    \hline
    \endfoot
    Pemrosesan pembayaran
        & Di luar \textit{scope}
        & Seluruh transaksi keuangan diproses SmartBank melalui API Gateway.
          PasarKita hanya mengirim \textit{payment request} dan membaca \textit{response}. \\
    \hline
    \textit{Routing} dan JWT validation (inter-service)
        & Di luar \textit{scope}
        & API Gateway (Kelompok 7) bertanggung jawab atas \textit{routing}, JWT
          \textit{issuance} antar-\textit{service}, dan \textit{logging}
          inter-\textit{service}. \\
    \hline
    Kalkulasi tarif dan \textit{tracking} pengiriman internal
        & Di luar \textit{scope}
        & LogistiKita (Kelompok 5) mengelola logika pengiriman. PasarKita hanya
          men-\textit{trigger} dan membaca \texttt{tracking\_id}. \\
    \hline
    Manajemen stok \textit{supplier}
        & Di luar \textit{scope}
        & SupplierHub (Kelompok 4) mengelola \textit{supply chain}. PasarKita
          hanya membaca stok lokal dari \textit{database} sendiri. \\
    \hline
    Laporan keuangan konsolidasi ekosistem
        & Di luar \textit{scope}
        & UMKM Insight (Kelompok 6) mengonsumsi \textit{ledger} dari SmartBank.
          PasarKita hanya menyediakan data transaksi lokal via \textit{analytics endpoint}. \\
    \hline
    Manajemen produk CRUD \textit{seller} (WarungPOS)
        & Di luar \textit{scope}
        & WarungPOS (Kelompok 3) memiliki katalog tersendiri. PasarKita mengelola
          katalog produk \textit{seller} yang terdaftar di platformnya sendiri. \\
    \hline
    Manajemen saldo dompet \textit{user}
        & Di dalam \textit{scope} (partial)
        & PasarKita menyimpan \textit{snapshot} saldo lokal dari SmartBank untuk
          keperluan validasi awal, namun \textit{authoritative balance} tetap di SmartBank. \\
    \hline
    Katalog produk, \textit{order}, dan ulasan
        & Di dalam \textit{scope}
        & PasarKita adalah \textit{system of record} untuk produk, \textit{order},
          \texttt{order\_items}, dan \textit{ratings} dalam ekosistemnya. \\
    \hline
\end{longtable}
